\par\medskip\noindent
\SGAInlineRefTarget{sga3:xix:lemma:1-11}\textbf{\textit{\SGARefLink{sga3:xix:lemma:1:11}{Lemma 1.11}.}}---
\textit{Let \(k\) be an algebraically closed field, \(G\) a smooth affine
connected \(k\)-group, and \(T\) a torus of \(G\).  Put
\[
  W(T)
  =
  \underline{\operatorname{Norm}}_{G}(T)/
  \underline{\operatorname{Centr}}_{G}(T),
\]
the Weyl group of \(G\) relative to \(T\).  The following conditions are
equivalent:
\begin{enumerate}
\item[(i)] \(G\) is reductive, \(T\) is maximal, and
  \(\operatorname{rgss}(G)=1\).
\item[(ii)] \(G\) is reductive, \(T\) is maximal, and \(G\ne T\); there
  exists a subtorus \(Q\) of \(T\), of codimension \(1\) in \(T\), that
  is central in \(G\).
\item[(iii)] \(G\) is not solvable and
  \(\dim(G)-\dim(T)\leq 2\).
\item[(iv)] \(W(T)\ne\{e\}\) and
  \(\dim(G)-\dim(T)\leq 2\).
\item[(v)] \(W(T)=(\mathbb Z/2\mathbb Z)_k\) and
  \(\dim(G)-\dim(T)=2\).
\end{enumerate}
Moreover, under these conditions there are exactly two roots of \(G\)
relative to \(T\), and they are opposite.  Under the conditions of
\textup{(ii)}, \(Q=\operatorname{rad}(G)\).}
\SGARefTarget{sga3:xix:lemma:1:11}

\SGAProof{Proof}
The implication (v)\(\Rightarrow\)(iv) is clear, and
(iv)\(\Rightarrow\)\SGARefLink{sga3:xix:lemma:1:11:item:iii}{(iii)} follows from \emph{Bible}, §6.1, Corollary 3 to
Theorem 1, or [Ch05], §6.2, Corollary 3 to Theorem 2.

Let us prove \SGARefLink{sga3:xix:lemma:1:11:item:iii}{(iii)}\(\Rightarrow\)\SGARefLink{sga3:xix:lemma:1:11:item:ii}{(ii)}.  Let \(U\) be the unipotent radical
of \(G\).  The group \(G/U\) is reductive and is not a torus (otherwise
\(G\) would be solvable); hence, by (\SGARefLink{sga3:xix:equation:1:10:4}{1.10.4}),
\[
  \dim(G/U)-\operatorname{rgred}(G/U)
  =
  \operatorname{Card}(R)
  \geq 2.
\]
On the other hand,
\[
  \operatorname{rgred}(G)
  =
  \operatorname{rgred}(G/U)+\operatorname{rgred}(U)
  =
  \operatorname{rgred}(G/U),
\]
and therefore
\[
\begin{aligned}
  \dim(G)-\dim(U)-\operatorname{rgred}(G)
    &=\dim(G/U)-\operatorname{rgred}(G/U)\\
    &\geq 2
     \geq \dim(G)-\dim(T)
     \geq \dim(G)-\operatorname{rgred}(G).
\end{aligned}
\]
It follows that \(\dim(U)=0\), so \(G\) is reductive, and that
\(\operatorname{rgred}(G)=\dim(T)\); thus \(T\) is maximal and
\(\dim(G)-\dim(T)=2\).  By \emph{Bible}, §10.4, Proposition 8, there
exists a singular torus \(Q\) of codimension \(1\) in \(T\).  Then
\(\underline{\operatorname{Centr}}_{G}(Q)\) is reductive (1.6.2(i)) and
not solvable (by the definition of a singular torus), and hence
\[
  \dim\bigl(\underline{\operatorname{Centr}}_{G}(Q)\bigr)
  -\operatorname{rgred}(G)\geq 2.
\]
This proves that
\(G=\underline{\operatorname{Centr}}_{G}(Q)\), and that \(Q\) is central
in \(G\).

Let us prove \SGARefLink{sga3:xix:lemma:1:11:item:ii}{(ii)}\(\Rightarrow\)\SGARefLink{sga3:xix:lemma:1:11:item:i}{(i)}.  We know that \(G/Q\) is reductive
by \SGARefLink{sga3:xix:paragraph:1-7}{1.7} and that \(\operatorname{rgred}(G/Q)=1\); hence
\(\operatorname{rgss}(G/Q)\) is either \(0\) or \(1\).  The first case is
impossible, since it would imply \(\operatorname{rgss}(G)=0\), and
therefore \(G=T\).  Thus
\[
  \operatorname{rgred}(G/Q)=\operatorname{rgss}(G/Q)=1,
\]
which proves that \(G/Q\) is semisimple.  Consequently \(Q\) is the
radical of \(G\) and \(\operatorname{rgss}(G)=1\).
